\chapter{Answers to the puzzles}
\label{app:answers}

Every answer here was run against the interpreter (puzzles 1--12) or loaded
into a real session (13--18). There are other good answers; if yours
works, it is right.

\section*{1. Temperatures}

\begin{scm}
(define (c->f c) (+ (/ (* c 9) 5) 32))
(define (frost? c) (<= c 0))
(list (c->f 100) (c->f 21) (frost? -2) (frost? 3))   ;=> (212 69.8 #t #f)
\end{scm}

Multiplying before dividing keeps \code{(c->f 100)} an exact \code{212}:
\code{/} returns an integer when the division comes out whole.

\section*{2. Seasons}

\begin{scm}
(define (season month)
  (cond ((member month '(12 1 2)) "winter")
        ((member month '(3 4 5)) "spring")
        ((member month '(6 7 8)) "summer")
        (else "autumn")))

(define (season month)
  (case month
    ((12 1 2) "winter")
    ((3 4 5) "spring")
    ((6 7 8) "summer")
    (else "autumn")))

(map season '(1 4 7 10 12))   ;=> ("winter" "spring" "summer" "autumn" "winter")
\end{scm}

\section*{3. The harvest}

\begin{scm}
(define harvest '(("tomato" . 12) ("basil" . 3) ("bean" . 40)))

(define (total h) (fold (lambda (sum p) (+ sum (cdr p))) 0 h))

(define (biggest h)
  (car (fold (lambda (best p) (if (> (cdr p) (cdr best)) p best))
             (car h) (cdr h))))

(define (total h)
  (let loop ((h h) (sum 0))
    (if (null? h) sum (loop (cdr h) (+ sum (cdar h))))))

(list (total harvest) (biggest harvest))   ;=> (55 "bean")
\end{scm}

\section*{4. Your own filter}

\begin{scm}
(define (keep pred l)
  (cond ((null? l) '())
        ((pred (car l)) (cons (car l) (keep pred (cdr l))))
        (else (keep pred (cdr l)))))

(define (keep pred l)
  (let loop ((l l) (acc '()))
    (cond ((null? l) (reverse acc))
          ((pred (car l)) (loop (cdr l) (cons (car l) acc)))
          (else (loop (cdr l) acc)))))

(keep odd? (iota 10))   ;=> (1 3 5 7 9)
\end{scm}

The accumulator is built by \code{cons}ing onto its front, so it comes out
backwards; one \code{reverse} at the end is cheaper than appending to the
end of a list on every step.

\section*{5. Zip}

\begin{scm}
(define (zip a b) (map cons a b))
(zip '(1 2 3) '(x y z))   ;=> ((1 . x) (2 . y) (3 . z))
\end{scm}

\section*{6. A label maker}

\begin{scm}
(define (label plant days)
  (if (= days 1)
      (format "~a: water every day" (string-upcase plant))
      (format "~a: water every ~a days" (string-upcase plant) days)))
(list (label "basil" 2) (label "mint" 1))
;=> ("BASIL: water every 2 days" "MINT: water every day")
\end{scm}

\section*{7. CSV to JSON}

\begin{scm}
(define (csv->json text)
  (let* ((rows (map (lambda (line) (string-split line ","))
                    (filter (lambda (l) (not (string-empty? l)))
                            (string-split text "\n"))))
         (keys (car rows)))
    (json-emit (map (lambda (row) (map cons keys row)) (cdr rows)))))

(csv->json "plant,water\nbasil,daily\ntomato,weekly")
;=> "[{\"plant\":\"basil\",\"water\":\"daily\"},{\"plant\":\"tomato\",\"water\":\"weekly\"}]"
\end{scm}

\code{(map cons keys row)} is puzzle 5's zip: it turns a row into an
association list with string keys, which \code{json-emit} writes as an
object.

\section*{8. A safe number parser}

\begin{scm}
(define (parse-ml text)
  (let ((words (string-split (string-trim text) " ")))
    (and (pair? words) (string->number (car words)))))

(define (must-parse-ml text)
  (or (parse-ml text) (error "not a quantity:" text)))

(parse-ml "250 ml")                     ;=> 250
(parse-ml "lots")                       ;=> #f
(guard (e (#t (format "~a ~s" (assoc-ref e 'message)
                                (car (assoc-ref e 'irritants)))))
  (must-parse-ml "lots"))               ;=> "not a quantity: \"lots\""
\end{scm}

\section*{9. A watering log}

\begin{scm}
(define (make-log)
  (let ((entries '()))
    (lambda args
      (if (null? args)
          (reverse entries)
          (set! entries (cons (car args) entries))))))

(define water-log (make-log))
(water-log "basil, Monday")
(water-log "mint, Tuesday")
(water-log)                 ;=> ("basil, Monday" "mint, Tuesday")
\end{scm}

\section*{10. Extend the calculator}

Two clauses before the general one, because \code{if} and \code{let} must
not evaluate all their parts first:

\begin{scm}
(define (calc expr env)
  (cond ((number? expr) expr)
        ((symbol? expr)
         (let ((b (assq expr env)))
           (if b (cdr b) (error "calc: unknown variable" expr))))
        ((and (pair? expr) (eq? (car expr) 'if))
         (if (not (= 0 (calc (cadr expr) env)))
             (calc (caddr expr) env)
             (calc (list-ref expr 3) env)))
        ((and (pair? expr) (eq? (car expr) 'let))
         (let* ((binding (car (cadr expr)))
                (value (calc (cadr binding) env)))
           (calc (caddr expr) (cons (cons (car binding) value) env))))
        ((pair? expr)
         (let ((op (car expr))
               (args (map (lambda (e) (calc e env)) (cdr expr))))
           (case op
             ((+) (apply + args)) ((-) (apply - args))
             ((*) (apply * args)) ((/) (apply / args))
             ((max) (apply max args))
             (else (error "calc: unknown operator" op)))))
        (else (error "calc: cannot evaluate" expr))))

(calc '(let ((x 4)) (if (- x 4) 1 2)) '())   ;=> 2
\end{scm}

Adding a binding by \code{cons}ing onto the front of \code{env} means an
inner \code{let} hides an outer one of the same name, because \code{assq}
finds the first match. That is lexical scope, in one line.

\section*{11. Hanoi, counted}

\begin{scm}
(define (hanoi-count n)
  (let loop ((k 0) (moves 0))
    (if (= k n) moves (loop (+ k 1) (+ (* 2 moves) 1)))))
(hanoi-count 10)   ;=> 1023
\end{scm}

Each extra disc doubles the work and adds one move: the recursion in
\code{hanoi}, read as arithmetic. (The closed form is $2^n - 1$.)

\section*{12. Roman numerals back}

\begin{scm}
(define (unroman s)
  (let loop ((s s) (table roman-table) (n 0))
    (cond ((string-empty? s) n)
          ((null? table) (error "not a roman numeral:" s))
          ((string-prefix? (cdar table) s)
           (loop (substring s (string-length (cdar table))) table
                 (+ n (caar table))))
          (else (loop s (cdr table) n)))))

(map unroman '("MMXXVI" "MCMXC" "XIV" "IV"))   ;=> (2026 1990 14 4)
\end{scm}

\section*{13. A harvest tool}

\begin{scm}
(define (harvest-rows)
  (map (lambda (line)
         (let ((cells (string-split line ",")))
           (cons (car cells) (string->number (string-trim (cadr cells))))))
       (filter (lambda (l) (not (string-empty? (string-trim l))))
               (string-split (fs-read "harvest.csv") "\n"))))

(register-tool "harvest-report"
  (lambda ()
    (let* ((rows (sort (harvest-rows) (lambda (a b) (> (cdr a) (cdr b)))))
           (total (fold (lambda (sum p) (+ sum (cdr p))) 0 rows)))
      (md-table '("crop" "count")
                (append (map (lambda (p) (list (car p) (cdr p))) rows)
                        (list (list "**total**" total))))))
  '((doc . "The harvest so far from harvest.csv, as a Markdown table, largest first, with a total.")
    (params . ())))
\end{scm}

\begin{txt}
| crop | count |
| --- | --- |
| bean | 40 |
| tomato | 12 |
| basil | 3 |
| **total** | 55 |
\end{txt}

\section*{14. Tests before commit}

\begin{scm}
(register-check "tests-before-commit" "deny"
  (lambda (prim args)
    (and (equal? prim "shell-exec")
         (pair? args)
         (equal? (car args) "git")
         (pair? (cdr args))
         (list? (cadr args))
         (member "commit" (cadr args)))))
\end{scm}

\begin{txt}
⚑ guard deny shell-exec: shell-exec is irreversible; user checker tests-before-commit
\end{txt}

Why it is weaker than it looks: it matches one spelling of one program.
\code{bash -c "git commit -m x"}, \code{/usr/bin/git commit}, a
\code{make} target or a registered tool that commits all get past it. A
check sees a call's arguments, not its intent. For a rule that matters,
deny the whole family (any \code{shell-exec} whose arguments mention
\code{commit} anywhere, say), and remember that every command no standing
approval covers is already asked about: of the judge first, by default,
and of you when it does not allow it.

\section*{15. Read before write}

\begin{scm}
(define read-paths '())

(define (call-path call)
  (let ((a (guard (e (else '())) (json-parse (assoc-ref call 'args "{}")))))
    (if (pair? a) (format "~a" (assoc-ref a "path" "")) "")))

(define rbw-base-run-calls run-calls)
(define (run-calls calls)
  (map (lambda (c)
         (let ((name (assoc-ref c 'name)) (path (call-path c)))
           (if (and (equal? name "fs-edit") (not (member path read-paths)))
               (skipped c "read the file before editing it")
               (let ((r (car (rbw-base-run-calls (list c)))))
                 (when (and (equal? name "fs-read") (assoc-ref r 'ok #f))
                   (set! read-paths (cons path read-paths)))
                 r))))
       calls))
\end{scm}

\code{skipped} is the helper from Chapter~\ref{ch:strategies}. The
\code{map} returns one result per call, in order: the one rule.

\section*{16. Second opinion}

\begin{scm}
(define-agent-tool "second-opinion"
  "Ask a second model the same question, without tools, for an independent opinion."
  '(("question" "string" "the question, with the context the other model needs"))
  (lambda (call args)
    (narrow-context "Answer the question briefly and say how sure you are." "second-opinion")
    (add-message (assoc-ref args "question"))
    (let ((reply (model-turn "local" '((tools . #f)))))
      (skip-tools (assoc-ref reply 'tool-calls '()) "tools are off: answer in text")
      (fold-context)
      (assoc-ref reply 'content ""))))
\end{scm}

The \code{skip-tools} matters: a local model offered no tools can still
answer with a call, and the context cannot fold while a call is
unanswered.

\section*{17. Plan only for big asks}

\begin{scm}
(define (big-ask? request)
  (or (> (string-length request) 200)
      (string-contains? (string-downcase request) "refactor")))

(define big-base-turn nodwyn-turn)
(define (nodwyn-turn request)
  (when (and (big-ask? request) (not (plan-recorded?)))
    (record-plan (list (cons 'decision "plan") (cons 'plan house-plan)
                       (cons 'reason "a big ask gets the house plan"))))
  (big-base-turn request))
\end{scm}

In a supervised run, a short request still gets a plan: the wrapped
\code{nodwyn-turn} underneath runs the default plan gate whenever
\code{(plan-required?)}, so the control model drafts one. And even if you
replaced \code{nodwyn-turn} entirely, the core would refuse anything
irreversible until a plan is recorded. You cannot plan less than the run
requires; you can only plan your own way.

\section*{18. Your own habit}

\begin{scm}
(define permhint-hints 2)
(define permhint-given 0)
(define permhint-hint
  "\n[hint: Permission denied: check whether the file is executable (ls -l), and run it with its interpreter or chmod +x it]")

(define permhint-base-run-calls run-calls)
(define (run-calls calls)
  (let ((results (permhint-base-run-calls calls)))
    (for-each
     (lambda (c r)
       (when (and (< permhint-given permhint-hints)
                  (equal? (assoc-ref c 'name) "shell-exec")
                  (string-contains? (format "~a" (assoc-ref r 'value "")) "Permission denied"))
         (set! permhint-given (+ permhint-given 1))
         (fold-result (assoc-ref c 'id)
                      (string-append (format "~a" (assoc-ref r 'value "")) permhint-hint))))
     calls results)
    results))
\end{scm}

\code{for-each} over two lists walks calls and results together, which
only works because \code{run-calls} keeps the one rule. The hint starts
with \code{"\textbackslash n[hint:"}, the convention \file{shell.scm}
recognises, so the shell habit still renders the result as plain text with
your hint after it.
